\scanpage{7}
\mspage{- 109 -}
% archive sheet number: 6
% lines 1–2
son idéal maximal, $X=\mathrm{Spec}(B)$ a deux points, $(0)$ et $\handwritten{\underline}{p}$, ce dernier étant le seul point fermé ; on a $\dim(X)=1$ mais $\dim\handwritten{\{}0\handwritten{\}}=0$.

% lines 3–6
\uline{Corollaire 2 .- Soit} $X$ \uline{un préschéma algébrique sur un corps $k$. Pou\handwritten{r} tout point} $x\handwritten{\in}X$, \uline{on a} $\dim_x(X)=\sup_i(\dim(X_i))$\handwritten{\unsure{,|$^\circ$}}, \uline{où} $X_i$ \uline{parcourt l'ense\handwritten{m}ble des composantes irréductibles de} $X$ \uline{contenant $x$. Si} $\handwritten{\{}x\handwritten{\}}$ \uline{est ferm\handwritten{é} on a} $\dim_x X=\operatorname{codim}_X x=\dim(\handwritten{\underline}{O}_x)$\handwritten{,} \inserted{\handwritten{et en tous cas $\dim_x X=\dim\underline O_{X,x}+\dim k(x)/k$}}

% lines 7–13
En effe\handwritten{t}, \deleted{pour tout ouvert $U$ contenant \deleted{\illegible{1}} $x$}, on a (n°1,cor.de la prop.1) $\dim_x(X)=\sup_i(\dim_x(X_i))$ et \deleted{\illegible{1}} d'après le cor.1 ; $\dim_x(X_i)=\allowbreak=\dim(X_i)$. D'autre part, comme \inserted{si $\handwritten{\{}x\handwritten{\}}$ est fermé,} toute chaîne de parties fermées irréductibles contenant $x$ est contenue dans\handwritten{\,}l'un des $X_i$, on a $\operatorname{codim}_X(x)\handwritten{=}\allowbreak=\sup_i(\operatorname{codim}_{X_i}(x))$, et en vertu de la formule (6) (applicable puisque $X_i$ est strictement équidimensionnel), $\operatorname{codim}_{X_i}(x)=\dim(X_i)$, d'où \deleted{\illegible{1}} la seconde assertion du corollaire.
\margin{\handwritten{$=\qquad?$}}

% lines 14–15
\uline{Corollaire 3 .- Soient} $X,Y$ \uline{deux préschémas algébriques sur un corps $k$. Pour tout $k$-morphisme} $f:X\handwritten{\to}Y$, \uline{on a} $\dim(\handwritten{\overline}{f(X)})\handwritten{\leq}\dim(X)$.
% slanted
\margin{\handwritten{\unsure{déjà vu en|a-t-on} \S{}1}}
\margin{\handwritten{\uline{égalité dans le cas quasi fini}}}

% lines 16–23
En considérant les composantes irréductibles de $\handwritten{\overline}{f(X)}$, on est aussitôt ramené au cas où \deleted{\illegible{2}} $Y$ est irréductible et $f$ dominant (n°\handwritten{1,} prop.1 (ii)) ; en outre, une au moins des composantes irréductible\handwritten{s} $X_i$ de $X$ est telle \inserted{alors} que \deleted{\illegible{1}} $f(X_i)$ soit dense \deleted{\illegible{1}} dans $Y$, donc on peu\handwritten{t} aussi supposer $X$ irréductible. Alors, si $x$ est le point générique de $X$, $f(x)=y$ est le point générique de $Y$, et comme il \deleted{\illegible{1}} correspond à $f$ un $k$-monomorphisme de $\handwritten{\underline{\kappa}}(y)$ dans $\handwritten{\underline{\kappa}}(x)$, on a \deleted{\illegible{1}} $\mathrm{deg.tr.}_k\,\handwritten{\underline{\kappa}}(y)\handwritten{\leq}\mathrm{deg.tr.}_k\,\handwritten{\underline{\kappa}}(x)$, d'où le corollaire.

% lines 24–29
Ici encore, ce résultat ne s'étend pas à des schémas intègres que\handwritten{l}conques. Par exemple, soient $B$ un anneau de valuation discrète, $K$ son corps des fractions, $X=\mathrm{Spec}(K)$, $Y=\mathrm{Spec}(B)$ ; à l'injection canonique $B\quad K$ correspond u\unsure{x|n} morphisme $f:X\handwritten{\to}Y$ \uline{de type fini} ($K$ étant en effet engendré par $1/\handwritten{\pi}$, si \deleted{\illegible{1}} $(\handwritten{\pi})$ est l'idéal maximal de $B$) ; mais on a \deleted{\illegible{1}} $\dim(Y)=1$ et $\dim(X)=0$.
% slanted
\margin{\handwritten{\deleted{\uline{Corollaire} $\dim_x X=\dim\underline O_x+\dim k(x)/k$}}}
\margin{\handwritten{déjà vu en \S{}1}}

% lines 30
\uline{Corollaire 4 .- \unsure{S}oient} $X,Y$ \uline{deux préschémas algébriques sur $k$ ; \handwritten{on a}}
